\section{Mejores prácticas de Claude Code}

Esta sección se basa en una publicación de Boris en X/Twitter que se volvió viral (99,000 marcadores), así como en su exposición detallada durante el programa.

\subsection{Flujo de trabajo con varios Claude en paralelo}

\begin{importantbox}{La forma de trabajar de Boris}
Boris ejecuta simultáneamente entre 5 y 10 instancias de Claude:
\begin{itemize}
    \item En la terminal local suele mantener varias pestañas abiertas, cada una con su propio checkout de git independiente
    \item Cuando se agotan las pestañas locales, continúa abriendo nuevas tareas en la versión web
    \item Cada mañana, al levantarse, inicia entre 2 y 3 tareas de Claude desde la aplicación de iOS
    \item Aproximadamente la mitad del código actual se produce a través de Claude en el teléfono
\end{itemize}
\end{importantbox}

Boris admite que hace un año jamás habría previsto que su forma de programar terminaría siendo así: repartiendo tareas entre el teléfono y la web, y recorriendo las distintas instancias para revisarlas, como quien administra un equipo de IA.

\subsection{Usar siempre el modelo más potente}

Boris recomienda usar siempre \textbf{Opus 4.5 with Thinking} (no Sonnet):

\begin{knowledgebox}{¿Por qué un modelo más grande resulta más económico?}
Se trata de un hallazgo contraintuitivo: aunque el precio por token de Opus 4.5 es más alto que el de Sonnet, al ser más inteligente:
\begin{itemize}
    \item Consume menos tokens en total para completar la misma tarea
    \item Requiere menos correcciones manuales
    \item El costo y el tiempo totales suelen ser \textbf{menores} que al usar un modelo más pequeño
\end{itemize}
Analogía: contratar a un ingeniero senior tiene un costo por hora más alto, pero como no necesita repetir el trabajo una y otra vez, el costo total puede ser menor que contratar a un ingeniero junior.
\end{knowledgebox}

\subsection{CLAUDE.md como base de conocimiento del equipo}

\begin{importantbox}{Mejores prácticas para CLAUDE.md}
\begin{itemize}
    \item El equipo comparte un archivo CLAUDE.md, integrado al repositorio de Git
    \item \textbf{Se actualiza varias veces por semana, entre todo el equipo}: en cuanto se detecta un error de Claude, se agrega de inmediato la regla de corrección a CLAUDE.md
    \item El formato es completamente libre: es solo un archivo de texto común, sin requisitos de sintaxis especiales
    \item Cada equipo mantiene su propio CLAUDE.md y es responsable de mantenerlo actualizado
\end{itemize}
\end{importantbox}

Boris mostró una captura de pantalla del CLAUDE.md del repositorio de Claude Code: su contenido es muy simple y sencillo, lo que demuestra la idea de que ``no se necesita un formato complejo, solo mantenimiento constante''.

\subsection{Plan Mode: planear antes de ejecutar}

\begin{importantbox}{``Once the plan is good, the code is good.''}
La mayoría de las sesiones de Boris comienzan en Plan Mode:
\begin{enumerate}
    \item Entra en Plan Mode y discute el plan con Claude una y otra vez
    \item Una vez confirmado el plan, cambia al modo Auto-accept Edits
    \item Claude por lo general puede completar todo el PR de una sola vez (one-shot)
\end{enumerate}
Esto coincide con la idea del Spec-driven Development, pero con más flexibilidad: el ``spec'' es solo un plan en forma de texto, sin necesidad de un formato específico.
\end{importantbox}

\subsection{GitHub Action e ingeniería compuesta}

La quinta recomendación de Boris tiene que ver con \textbf{acumular de manera continua el conocimiento del equipo durante la revisión de código}:

\begin{enumerate}
    \item Instalar la GitHub App de Claude en el repositorio mediante \texttt{/install-github-action}
    \item En la revisión del PR, mencionar a @claude para que aplique modificaciones o actualice CLAUDE.md
    \item Esto garantiza que un mismo tipo de problema \textbf{nunca necesite señalarse dos veces}
\end{enumerate}

\begin{knowledgebox}{Ingeniería compuesta (Compound Engineering)}
Boris cita el concepto de ``ingeniería compuesta'' de Dan Schipper: cuando trabajaba en Meta, Boris mantenía una hoja de cálculo con los problemas que surgían en las revisiones de código, y cuando un problema se repetía entre 5 y 10 veces, escribía una regla de lint para automatizar su detección. Ahora, CLAUDE.md más GitHub Action es la versión de esta idea para la era de la IA: en cuanto se detecta un problema, se le pide de inmediato a Claude que actualice la base de conocimiento, para garantizar que el mismo error no se repita.
\end{knowledgebox}

\subsection{Darle a Claude la capacidad de verificar su propia salida}

Boris considera que esto es uno de los \textbf{tres factores clave} para mejorar el desempeño de Claude Code (junto con usar Opus y mantener CLAUDE.md):

\begin{warningbox}{La analogía del ``pintor con los ojos vendados''}
Si un pintor tuviera que trabajar con los ojos vendados, sin importar su habilidad, no podría producir una buena obra. De la misma manera, si Claude no puede ver el resultado de ejecutar el código que escribió, la calidad de su salida se ve muy afectada. Por eso:
\begin{itemize}
    \item Al construir una aplicación web, es indispensable habilitar la extensión de Chrome para que Claude pueda ver el resultado en el navegador
    \item Siempre hay que dejar que Claude ejecute las pruebas, inicie el servidor y revise la salida del simulador
    \item Permitir que Claude se verifique y se corrija a sí mismo es el paso decisivo para pasar de ``aceptable'' a ``excelente''
\end{itemize}
\end{warningbox}

\subsection{Resumen de la sección}

Las mejores prácticas de Claude Code según Boris se pueden resumir en cinco principios centrales:
\begin{enumerate}
    \item \textbf{Paralelismo}: ejecutar varios Claude a la vez, asignando tareas como quien administra un equipo
    \item \textbf{El modelo más potente}: usar siempre Opus 4.5 with Thinking
    \item \textbf{Acumulación de conocimiento}: mantener CLAUDE.md de forma continua como base de conocimiento del equipo
    \item \textbf{Planear primero}: usar Plan Mode para asegurar que el plan sea correcto antes de pasar a la ejecución automática
    \item \textbf{Autoverificación}: permitir que Claude vea y verifique su propia salida
\end{enumerate}
